\section{計算法の詳細}\label{app:computation}
\paragraph{評価と補間．}
$\overline{\rCS}_\lambda$ と $\hf_\rho$ は，31 ビットの素数 $p$ を法として格子の点で評価し，$(A,Q,T)$ の各変数について補間して求める．
係数が $p/2$ を超えるときは 2〜3 個の素数で補間し，中国剰余定理（Garner の方法）で復元する．
得られた多項式は，補間に使っていない素数の点で検証する．
評価の点はまとまりごとに保存し，計算を途中から再開できるようにした．
\paragraph{genus 展開．}
補題~\ref{lem:G} と同じく $v=\bq^2\bt$，$w=v+v^{-1}$，$w_0=\bt+\bt^{-1}$ とおくと $G=w-w_0$ である．各行を $w$ の多項式に直し，Horner 法で $G$ の冪に展開する．
\paragraph{計算時間．}
計算は Apple M3 Pro（メモリ 36 GB）と Mac mini M1（メモリ 8 GB）で行った．
最も重いのは $K10$ である．$\rCS_{\yng(3)}$ は 2 つの素数で約 4 時間（14440 秒），$\rCS_{\yng(2,1)}$ は約 2 時間（7387 秒）かかった．
3 箱の $\wh f_\rho$ は 3 つの素数を要し，4050 秒，最大メモリ 17.6 GB であった．
$K7$ と $K8$ の $|\rho|\le3$ は，それぞれ 12275 秒と 22853 秒で，最大メモリは 2 GB 未満であった．
\paragraph{実装の検証．}
実装は，小さな入力で各段を確かめるテストで検証した．
対象は，KN の値の再現，Cherednik--Danilenko の明示例，ブレイドから計算した HOMFLY-PT 多項式，補間と中国剰余定理，genus 展開の二つの実装の一致などである．

\section{正規化の導出}\label{app:normalization}
\paragraph{(A) ケーブル．}
framing 演算子 $\tau:P_\mu\mapsto T_\mu P_\mu$ について，$\tau^jP_\lambda[p_k\mapsto P_{k(m,n)}]\cdot1=P_\lambda[p_k\mapsto P_{k(m,n+jm)}]\cdot1$ が成り立つ（計算結果~\ref{cr:tau}）．
ケーブルの構成は companion について線形である．companion に (N2$'$) を適用すると，パターンの傾きが $(m,-s')$ から $(m,-n)$（Seifert framing）にずれる．
仮説~\ref{hyp:sat} と合わせると，$M_{C_{m,n}K}=M_K^mM_{T_{m,n}}$ と $F=mn$ が従う．
\paragraph{(B) トーラス結び目．}
(I) Gorsky--Neguţ の正規化前の reduced 値の $a^0$ の係数は $(-t/q)^{m-1}C_{m,n}(1/q,t)$ である（$C_{m,n}$ は有理 $q,t$-Catalan 数）．
これは Gorsky--Neguţ の予想であり，Bergeron--Garsia--Leven--Xin の $Q_{m,n}$ を介して Mellit~\cite{Mellit16} の定理から従う．
(II) 最低 $a$ 次数 $\delta$ は Franks--Williams の正のブレイドについての定理から従う．
(III) $t=q$ での符号と $q$ の冪は，$C_{m,n}$ の $q,t$ 対称性から決まる．
(IV) $(q/t)$ の指数は，Mellit~\cite{Mellit17} の Dyck 路の表示と計算結果~\ref{cr:mellit} の式から決まる．
以上から $M_{T_{m,n}}=(-q/t)^{m-1}(a(q/t)^{1/2})^{\delta}$ である．

\section{非代数的なケーブル}\label{app:nonalg}
本論文で扱う非代数的なケーブルは，$T_{2,3}$ の $(2,1)$-，$(2,3)$-，$(2,5)$-，$(2,7)$-，$(2,-1)$-ケーブル（順に $K2,K4,K5a,K5b,K5c$）と，$T_{2,5}$ の $(2,1)$-ケーブル（$K6b$）である．
非代数的なケーブル（$n_2<m_1n_1m_2$）では，DAHA の値はホモロジーの Poincaré 多項式ではない（スペクトル系列の必要条件を満たさない）．
したがって KN の正値性の前提を満たさない．表~\ref{tab:nonalg} に判定を示す．
\begin{table}[h]
\centering
\caption{非代数的なケーブルの判定（framing 0）．各欄は「spin 表示の負の項の数/項数; genus 表示の負の項の数/項数」（全体の符号はそれぞれで良い方）．}\label{tab:nonalg}
\resizebox{\textwidth}{!}{% scripts/paper_summary.py が生成。\checkmark(n) は G 展開の前後とも負の項 0（n は前の項数）。それ以外は 前の負/前の項数; 後の負/後の項数（前後それぞれで全体の符号は良い方）。framing 0。
\begin{tabular}{llcccccc}
\hline
結び目 & 分類 & $\Y{1}$ & $\Y{2}$ & $\Y{1,1}$ & $\Y{3}$ & $\Y{2,1}$ & $\Y{1,1,1}$ \\
\hline
K2 & non & 10/47; 6/48 & 152/1136; 41/1045 & 174/1273; 51/1190 & 967/8525; 253/7852 & 1078/9527; 292/8774 & 1156/9990; 318/9254 \\
K4 & non & 18/42; 20/43 & 286/1109; 195/952 & 330/1255; 224/1090 & 2561/8614; 1749/7420 & 2850/9646; 1989/8325 & 3025/10121; 2128/8795 \\
K5a & non & 8/38; 5/47 & 409/1226; 332/1066 & 460/1377; 376/1210 & 4441/9758; 3631/8278 & 4982/10844; 4066/9235 & 5164/11348; 4263/9733 \\
K5b & non & \checkmark\,(48) & 136/1395; 66/1305 & 149/1537; 78/1452 & 2097/10706; 1060/9612 & 2358/11707; 1181/10553 & 2351/12202; 1228/11040 \\
K5c & non & 5/57; 0/69 & 78/1307; 0/1299 & 77/1443; 0/1433 & 400/9501; 6/9123 & 423/10494; 7/10053 & 376/10864; 0/10444 \\
K6b & non & 30/133; 10/178 & 576/4400; 81/4367 & 618/4677; 91/4656 & 3626/33214; 512/32175 & 3766/35136; 541/34032 & 3878/36067; 560/34980 \\
\hline
\end{tabular}
}
\end{table}
破れ方は三つの型に分かれる．
(I) $T_{2,3}$ の $(2,3)$，$(2,5)$-ケーブル（$K4$，$K5a$）では，色付きでない HOMFLY-PT の段階で，偶奇のそろった refinement による正値性が不可能である．
(II) $(2,1)$-ケーブル（$K2$，$K6b$）と $(2,-1)$-ケーブル（$K5c$）では，HOMFLY-PT は正値性と両立するが，DAHA の $\rCS_\square$ 自体が負の係数をもつ．
(III) $T_{2,3}$ の $(2,7)$-ケーブル（$K5b$）の色付きの場合の破れは framing に依存する．偶数の framing $x$ を調べると，2 箱は $x\le-12$，3 箱は $x\le-14$ で前後とも正になる（$|x|\le24$ の偶数で計算）．
$K5c$ は展開の後だけが（ほぼ）正になる例で，補題~\ref{lem:G} の逆が成り立たないことを示す．

\section{データとコード}\label{app:data}
$\overline{\rCS}_\lambda$，$\hf_\rho$，spin 表示と genus 表示の全ての値と，計算・判定のコード，テストを公開する．
